% ============================================================
% Verbatim prompt appendix, generated directly from the released code.
% Every string below is copied character-for-character from the scripts
% listed in prompt_sources.md; templates are shown with {placeholders}
% exactly as they appear in the code, followed by one fully instantiated
% example where construction is programmatic.
% ============================================================

\section*{Verbatim prompt templates}

API experiments use the OpenAI chat format; local-checkpoint experiments apply each model's own chat template (noted per section). The non-Likert benchmark generalizes the A.2/A.3 templates to $k$-option items with a shortened scale-anchored phrasing; its exact builders are in \texttt{silicon/nonlikert.py} (see the source map). Prompts below are annotated
with their role (\texttt{system} / \texttt{user} / \texttt{assistant}).
Where a prompt is assembled programmatically we first show the template
with \texttt{\{placeholders\}} and then one fully instantiated example.
Prompt text is reproduced character-for-character, with one exception:
prompts that are a single long line (or one continuous paragraph) in
the code are soft-wrapped here to fit the page width; explicit
\verb|\n| line breaks in the real prompts are shown as actual line
breaks.

% ------------------------------------------------------------
\subsection*{A.1\quad Demographic persona template (Argyle-style silicon sampling)}

The one-sentence persona template shared, via import, by every
experiment that samples personas (\texttt{silicon/opinionqa.py}
\texttt{sample\_personas}). Personas are drawn from the Pew ATP W92 joint
demographic distribution (seed 42) and rendered with this template:

\begin{verbatim}
You are a {race} {gender} {age} {educ}, {inc}, who is {party}.
\end{verbatim}

The slot fillers are fixed mappings from Pew category labels:

\begin{verbatim}
AGE:    "18-29" -> "in their 20s"
        "30-49" -> "in their 30s or 40s"
        "50-64" -> "in their 50s or early 60s"
        "65+"   -> "in their late 60s or older"
EDUC:   "H.S. graduate or less" -> "with a high-school education"
        "Some College"          -> "with some college experience"
        "College graduate+"     -> "with a college degree or more"
RACE:   "White non-Hispanic" -> "non-Hispanic white"
        "Black non-Hispanic" -> "non-Hispanic Black"
        "Hispanic"           -> "Hispanic"
        "Asian non-Hispanic" -> "non-Hispanic Asian American"
        "Other"              -> "of another racial background"
PARTY:  "Dem/Lean Dem"       -> "Democrat or Democratic-leaning"
        "Rep/Lean Rep"       -> "Republican or Republican-leaning"
        "DK/Refused/No lean" -> "politically independent"
SEX:    "Female" -> "woman"   "Male" -> "man"
        "In some other way" -> "person"
INCOME: "Lower income"  -> "in a lower-income household"
        "Middle income" -> "in a middle-income household"
        "Upper income"  -> "in an upper-income household"
\end{verbatim}

Fully instantiated example:

\begin{verbatim}
You are a non-Hispanic white woman in their 30s or 40s with a college
degree or more, in a middle-income household, who is Democrat or
Democratic-leaning.
\end{verbatim}

% ------------------------------------------------------------
\subsection*{A.2\quad Argyle baseline prompt, the 100-item expansion run (per-item results)}

Used for the 100-item baseline (100 personas $\times$ 1 call each,
$T{=}1.0$). The user message template (\texttt{build\_argyle}) is:

\begin{verbatim}
{phrasing}

Options:
1. {label shown at position 1}
2. {label shown at position 2}
3. {label shown at position 3}
4. {label shown at position 4}
5. {label shown at position 5}

Respond with just the number (1-5).
\end{verbatim}

where \texttt{\{phrasing\}} is the item question under the selected
phrasing template (A.3) and the five displayed labels are the item's
options permuted by the selected ordering (A.3). The baseline condition
fixes ordering = \texttt{ascending}, phrasing = \texttt{direct},
position = \texttt{system\_then\_user}. Fully instantiated baseline
example (item \texttt{AGEMAJ\_W41}):

\begin{verbatim}
[system]
You are a non-Hispanic white woman in their 30s or 40s with a college
degree or more, in a middle-income household, who is Democrat or
Democratic-leaning.

[user]
According to the U.S. Census Bureau, by the year 2050, the number of
people in the population who are 65 or older will outnumber people
younger than 18. In terms of its impact on the country, do you t

Options:
1. A very good thing
2. A somewhat good thing
3. A somewhat bad thing
4. A very bad thing
5. Neither a good nor bad thing

Respond with just the number (1-5).
\end{verbatim}

(The question text is stored verbatim in \texttt{item\_list\_v2.json};
the truncation of this particular item's stem is present in the data
file and therefore in the prompt.)

% ------------------------------------------------------------
\subsection*{A.3\quad PPA perturbation grid (prompt-perturbation averaging)}

The PPA condition re-uses the A.2 template but randomizes, per call,
one level from each of three factors (random seed 42 + item-key hash).

\textbf{Factor 1 — option orderings} (4 levels). Each sequence lists,
for displayed positions 1--5, which canonical option index is shown
there; responses are mapped back to canonical indices before analysis:

\begin{verbatim}
ascending:  [1, 2, 3, 4, 5]
descending: [5, 4, 3, 2, 1]
shuffled_a: [3, 5, 1, 4, 2]
shuffled_b: [2, 4, 1, 5, 3]
\end{verbatim}

\textbf{Factor 2 — question phrasings} (3 levels), applied to the raw
question text \texttt{\{q\}}:

\begin{verbatim}
direct:       {q}
indirect:     On a scale from 1 to 5, how would you rate the following
              question: {q}
third_person: How would you, in your role, respond to this question: {q}
\end{verbatim}

(the \texttt{indirect} and \texttt{third\_person} templates are single
lines in the code; they are wrapped here only for the page width).

\textbf{Factor 3 — persona position} (3 levels), controlling where the
persona text \texttt{p} goes relative to the A.2 user message
\texttt{user}:

\begin{verbatim}
system_then_user: [{role: system, content: p},
                   {role: user,   content: user}]

user_only:        [{role: user,   content: p + "\n\n" + user}]

user_with_recap:  [{role: system, content: p},
                   {role: user,   content:
                      "Given this background, answer:\n\n" + user}]
\end{verbatim}

% ------------------------------------------------------------
\subsection*{A.4\quad Describe-pathway prompt, the 100-item expansion run}

Single call per item, $T{=}0$, JSON response format enforced. User
message (\texttt{build\_describe}; no system prompt):

\begin{verbatim}
Estimate the distribution of responses to the following Pew survey
item across a representative sample of U.S. adults. Output a JSON
object with a single key 'distribution' whose value is a list of
5 probabilities summing to 1.0, in the order of options 1 through 5.

Question: {question}

Options:
  1. {option 1}
  2. {option 2}
  3. {option 3}
  4. {option 4}
  5. {option 5}
\end{verbatim}

(In the source this is one continuous paragraph followed by the
question and options; the sentence text is verbatim.)

% ------------------------------------------------------------
\subsection*{A.5\quad Verbalized-sampling (VS) prompt (the 100-item expansion run and the earlier VS robustness run)}

Single call per item, $T{=}0$, JSON response format enforced. User
message (\texttt{build\_vs}; no system prompt):

\begin{verbatim}
For the following survey question, list each possible response option
along with the proportion of US adults who would select it. Output as
JSON: a single key 'responses' whose value is a list of objects, each
with 'response' (string) and 'probability' (number). Probabilities
must sum to 1. Be precise; this is being used for downstream
estimation.

Question: {question}

Response options:
  - {option 1}
  - {option 2}
  - {option 3}
  - {option 4}
  - {option 5}
\end{verbatim}

% ------------------------------------------------------------
\subsection*{A.6\quad Demographically matched Argyle persona (matched-Argyle protocol, paper Appendix on matched Argyle)}

The matched-Argyle protocol replaces the marginal-resampled personas
of A.1 with personas built from actual Pew respondent rows for each
item (sampled with survey weights, seed 42 + item-key hash). The persona template is
byte-identical to A.1:

\begin{verbatim}
You are a {race} {gender} {age} {educ}, {inc}, who is {party}.
\end{verbatim}

with the same category mappings as A.1 (fallbacks: \texttt{person} for
unmapped SEX, raw Pew label for unmapped AGE/RACE/EDUC, empty string
for unmapped PARTY/INCOME). Each matched persona answers the fixed
(ascending-order, direct-phrasing) prompt:

\begin{verbatim}
[system]
{matched persona descrip}

[user]
{question}

Options:
1. {option 1}
2. {option 2}
3. {option 3}
4. {option 4}
5. {option 5}

Respond with just the number (1-5).
\end{verbatim}

% ------------------------------------------------------------
\subsection*{A.7\quad Cross-family RNG probe prompts}

Five conditions, 200 calls each at $T{=}1.0$, user message only:

\begin{verbatim}
R1_int100: Pick a random integer uniformly between 1 and 100. Output
           only the integer.
R2_int10:  Pick a random integer uniformly between 1 and 10. Output
           only the integer.
R3_float:  Output a sample from the uniform distribution on the
           interval [0, 1]. Output only the number.
R4_coin:   Flip a fair coin. Output H or T.
R5_letter: Pick a random letter from the English alphabet uniformly at
           random. Output only the letter.
\end{verbatim}

(each prompt is a single line in the code).

% ------------------------------------------------------------
\subsection*{A.8\quad Cross-family seven-target panel prompts}

Seven target distributions, 200 calls each at $T{=}1.0$, user message
only:

\begin{verbatim}
point_5:       Always output the digit 5. Reply with just the digit.
mixture_54:    Output digit 5 with probability 70% and digit 4 with
               probability 30%. Reply with just the digit.
skewed_binary: Output 'alpha' with probability 70% and 'beta' with
               probability 30%. Reply with just the word.
bimodal_15:    Output digit 1 about 50% and digit 5 about 50%, never
               2/3/4. Reply with just the digit.
three_way:     Output A about 50%, B about 30%, C about 20%. Reply with
               just the letter.
skewed_5way:   Output a digit: 1 about 40%, 2 about 30%, 3 about 15%,
               4 about 10%, 5 about 5%. Reply with just the digit.
uniform_15:    Pick a random digit from 1 to 5 uniformly at random.
               Reply with just the digit.
\end{verbatim}

(each prompt is a single line in the code).

% ------------------------------------------------------------
\subsection*{A.9\quad KNOWS/DOES task prompts, V0 direct sampling (cross-family KNOWS/DOES panel)}

Five tasks; V0 issues the bare instruction 200 times at $T{=}1.0$:

\begin{verbatim}
uniform_digit: Pick a random digit from 1 to 5 uniformly at random.
               Reply with just the digit.
fair_coin:     Flip a fair coin. Output H or T.
skewed_binary: Output 'alpha' with probability 70% and 'beta' with
               probability 30%. Reply with just the word.
bimodal_5way:  Output a digit from 1 to 5 where 1 appears 40%, 2
               appears 10%, 3 appears 0%, 4 appears 10%, 5 appears 40%.
               Reply with just the digit.
skewed_5way:   Output a digit from 1 to 5 where 1 appears 50%, 2
               appears 20%, 3 appears 10%, 4 appears 10%, 5 appears
               10%. Reply with just the digit.
\end{verbatim}

(each prompt is a single line in the code).

% ------------------------------------------------------------
\subsection*{A.10\quad V7 algorithmic chain-of-thought (inverse-CDF) prompt (cross-family KNOWS/DOES panel)}

System-prompt template (\texttt{build\_v7\_messages}); the user message
is the task's V0 prompt from A.9. \texttt{\{cum\_str\}} enumerates the
cumulative distribution, \texttt{\{sup\_str\}} the support:

\begin{verbatim}
[system]
Follow this inverse-CDF procedure and output exactly three lines:
Line 1: 'F={cum_str}'
Line 2: 'u=<random number in [0,1]>'
Line 3: 'answer=<your choice from {{sup_str}}>'

[user]
{v0_prompt}
\end{verbatim}

Fully instantiated example (task \texttt{uniform\_digit}):

\begin{verbatim}
[system]
Follow this inverse-CDF procedure and output exactly three lines:
Line 1: 'F=F(1)=0.20, F(2)=0.40, F(3)=0.60, F(4)=0.80, F(5)=1.00'
Line 2: 'u=<random number in [0,1]>'
Line 3: 'answer=<your choice from {1, 2, 3, 4, 5}>'

[user]
Pick a random digit from 1 to 5 uniformly at random. Reply with just
the digit.
\end{verbatim}

% ------------------------------------------------------------
\subsection*{A.11\quad V8 list-of-N prompt (cross-family KNOWS/DOES panel)}

User-message template (\texttt{build\_v8\_messages}); one call yields
20 samples, 10 calls per cell. \texttt{\{freq\_str\}} enumerates
expected counts \texttt{round(p\_i * 20)} per support element:

\begin{verbatim}
Generate a list of 20 answers where {freq_str}. Output as JSON with
key 'samples' whose value is a list of 20 strings.
\end{verbatim}

Fully instantiated example (task \texttt{uniform\_digit}):

\begin{verbatim}
Generate a list of 20 answers where '1' about 4 times, '2' about 4
times, '3' about 4 times, '4' about 4 times, '5' about 4 times. Output
as JSON with key 'samples' whose value is a list of 20 strings.
\end{verbatim}

% ------------------------------------------------------------
\subsection*{A.12\quad V9 parametric-specification prompt (cross-family KNOWS/DOES panel)}

User-message template (\texttt{build\_v9\_messages}); one call at
$T{=}0$, then 200 external NumPy draws (seed 42) from the returned
distribution:

\begin{verbatim}
Output the probability distribution over {{sup_str}} that matches:
{prob_str}. Output as JSON with key 'distribution' whose value is a
list of {K} probabilities summing to 1.0.
\end{verbatim}

Fully instantiated example (task \texttt{uniform\_digit}):

\begin{verbatim}
Output the probability distribution over {1, 2, 3, 4, 5} that matches:
P(1)=0.20, P(2)=0.20, P(3)=0.20, P(4)=0.20, P(5)=0.20. Output as JSON
with key 'distribution' whose value is a list of 5 probabilities
summing to 1.0.
\end{verbatim}

% ------------------------------------------------------------
\subsection*{A.13\quad Correction-strategy prompts V1--V6 (the correction-strategy experiments and the extended non-uniform correction runs)}

Six correction system prompts, each paired with the task's base user
prompt (V0). The base tasks in the correction-strategy experiments are:

\begin{verbatim}
uniform_digit V0: Pick a random digit from 1 to 5 uniformly at random.
                  Reply with just the digit.
coin_flip V0:     Flip a fair coin. Report heads or tails. Reply with
                  just the word.
\end{verbatim}

with task parameters (verbatim from the task-parameter file in the
supplement):

\begin{verbatim}
uniform_digit:
  baseline_behavior_str: digit 1: 0.5%, digit 2: 0.5%, digit 3: 94.5%,
                         digit 4: 5%, digit 5: 0%
  mode_label: digit 3        mode_pct: 94.5%
  target_phrase: uniformly-random digit so that across many responses
                 you give each digit 1-5 about 20% of the time
  target_outcome_short: uniform random digit (each digit 1-5 about 20%
                        of the time)

coin_flip:
  baseline_behavior_str: heads: 97.5%, tails: 2.5%
  mode_label: heads          mode_pct: 97.5%
  target_phrase: fair coin outcome so that across many responses you
                 give heads about 50% of the time and tails about 50%
                 of the time
  target_outcome_short: fair coin outcome (heads about 50%, tails
                        about 50%)
\end{verbatim}

\textbf{V1 verbatim feedback} (system-prompt template):

\begin{verbatim}
I have measured your actual behavior on the prompt
"{base_user_prompt}". Across 200 runs, you produced:
{baseline_behavior_str}. This is far from the target. For this task,
please correct: produce a {target_phrase}.
\end{verbatim}

Fully instantiated (uniform\_digit):

\begin{verbatim}
I have measured your actual behavior on the prompt "Pick a random
digit from 1 to 5 uniformly at random. Reply with just the digit.".
Across 200 runs, you produced: digit 1: 0.5%, digit 2: 0.5%, digit 3:
94.5%, digit 4: 5%, digit 5: 0%. This is far from the target. For this
task, please correct: produce a uniformly-random digit so that across
many responses you give each digit 1-5 about 20% of the time.
\end{verbatim}

\textbf{V2 apologetic framing}:

\begin{verbatim}
I'm sorry for the confusion in our previous conversation. I realize my
earlier instructions may not have been clear. Let me try again with
clearer guidance. When you respond to the prompt "{base_user_prompt}",
I would like you to produce a truly {target_outcome_short}.
\end{verbatim}

Fully instantiated (uniform\_digit):

\begin{verbatim}
I'm sorry for the confusion in our previous conversation. I realize my
earlier instructions may not have been clear. Let me try again with
clearer guidance. When you respond to the prompt "Pick a random digit
from 1 to 5 uniformly at random. Reply with just the digit.", I would
like you to produce a truly uniform random digit (each digit 1-5 about
20% of the time).
\end{verbatim}

\textbf{V3 third-person framing}:

\begin{verbatim}
The assistant has been producing {mode_label} in approximately
{mode_pct} of trials on the prompt "{base_user_prompt}". The target
behavior is a {target_outcome_short}. The assistant should now produce
outputs consistent with this target.
\end{verbatim}

Fully instantiated (uniform\_digit):

\begin{verbatim}
The assistant has been producing digit 3 in approximately 94.5% of
trials on the prompt "Pick a random digit from 1 to 5 uniformly at
random. Reply with just the digit.". The target behavior is a uniform
random digit (each digit 1-5 about 20% of the time). The assistant
should now produce outputs consistent with this target.
\end{verbatim}

\textbf{V4 explicit bias naming}:

\begin{verbatim}
You have a known bias toward producing {mode_label} when asked to
respond to "{base_user_prompt}". This is a documented failure mode in
language models. Please consciously overcome this bias on the
following task and produce a {target_phrase}.
\end{verbatim}

Fully instantiated (uniform\_digit):

\begin{verbatim}
You have a known bias toward producing digit 3 when asked to respond
to "Pick a random digit from 1 to 5 uniformly at random. Reply with
just the digit.". This is a documented failure mode in language
models. Please consciously overcome this bias on the following task
and produce a uniformly-random digit so that across many responses you
give each digit 1-5 about 20% of the time.
\end{verbatim}

\textbf{V5 in-context demonstration (20 shots)}. The system prompt is
task-specific in the correction-strategy experiments:

\begin{verbatim}
uniform_digit:
Below are examples of correctly producing uniformly-random digits from
1 to 5. Each subsequent response should be drawn from the same uniform
distribution.

coin_flip:
Below are examples of correctly producing fair coin flips. Each
subsequent response should be drawn from the same fair-coin
distribution.
\end{verbatim}

The message list then contains 20 in-context turns: for each
demonstration value the pair
(\texttt{user}: base user prompt, \texttt{assistant}: value),
followed by a final \texttt{user}: base user prompt. The verbatim
demonstration sequences are:

\begin{verbatim}
uniform_digit: 4, 1, 5, 2, 3, 4, 5, 1, 2, 3, 5, 1, 4, 2, 3, 4, 1, 5,
               3, 2
coin_flip:     heads, tails, heads, tails, tails, heads, tails, heads,
               heads, tails, tails, heads, tails, heads, heads, tails,
               tails, heads, tails, heads
\end{verbatim}

\textbf{V6 step-by-step reasoning}. Task-specific system prompts in
the correction-strategy experiments:

\begin{verbatim}
uniform_digit:
When you respond to the next prompt, follow these steps:
Step 1: Acknowledge that uniform random sampling means each outcome
has equal probability.
Step 2: Identify what the support of the distribution is (digits 1, 2,
3, 4, 5).
Step 3: Mentally simulate a process that produces each outcome with
probability 1/5.
Step 4: Report the outcome of that simulation as your response.

coin_flip:
When you respond to the next prompt, follow these steps:
Step 1: Acknowledge that a fair coin has two outcomes with equal
probability.
Step 2: Identify the support (heads, tails).
Step 3: Mentally simulate a process that produces each outcome with
probability 1/2.
Step 4: Report the outcome of that simulation as your response.
\end{verbatim}

(each \texttt{Step} is a single line in the code). For V6 the user
message is the base user prompt with the verbatim suffix:

\begin{verbatim}
 Apply the steps above silently and respond with only the result.
\end{verbatim}

\textbf{Extended non-uniform correction runs (three non-uniform
tasks).} The same V1--V4
templates are instantiated with these task parameters (verbatim):

\begin{verbatim}
skewed_binary:
  V0: Output 'alpha' with probability 70% and 'beta' with probability
      30%. Output one of: alpha, beta. Reply with just the word.
  baseline_behavior_str: alpha: 100%, beta: 0%
  mode_label: alpha          mode_pct: 100%
  target_phrase: [alpha, beta]-valued response so that across many
                 responses you give alpha about 70% of the time and
                 beta about 30% of the time
  target_outcome_short: skewed binary (alpha about 70%, beta about 30%)
  icl_examples: alpha, beta, alpha, alpha, beta, alpha, alpha, alpha,
                beta, alpha, alpha, beta, alpha, alpha, alpha, beta,
                alpha, alpha, alpha, beta

bimodal_5way:
  V0: Output a digit from 1 to 5 where 1 appears 40% of the time, 2
      appears 10%, 3 appears 0%, 4 appears 10%, and 5 appears 40%.
      Reply with just the digit.
  baseline_behavior_str: digit 1: 100%, others: 0%
  mode_label: digit 1        mode_pct: 100%
  target_phrase: digit drawn from [P(1)=0.40, P(2)=0.10, P(3)=0.00,
                 P(4)=0.10, P(5)=0.40]
  target_outcome_short: bimodal on {1,2,4,5} with peaks at 1 and 5
                        (40% each); 3 is impossible
  icl_examples: 1, 5, 4, 1, 5, 1, 5, 2, 1, 5, 1, 5, 4, 5, 1, 5, 5, 1,
                2, 5

skewed_5way:
  V0: Output a digit from 1 to 5 where 1 appears 50%, 2 appears 20%,
      3 appears 10%, 4 appears 10%, 5 appears 10%. Reply with just the
      digit.
  baseline_behavior_str: digit 1: 100%, others: 0%
  mode_label: digit 1        mode_pct: 100%
  target_phrase: digit drawn from [P(1)=0.50, P(2)=0.20, P(3)=0.10,
                 P(4)=0.10, P(5)=0.10]
  target_outcome_short: skewed 5-way (1 dominant at 50%, others
                        20/10/10/10)
  icl_examples: 1, 1, 2, 1, 3, 1, 2, 4, 1, 1, 5, 2, 1, 1, 3, 2, 1, 4,
                1, 5
\end{verbatim}

In the extended non-uniform correction runs the V5 system prompt is
the generic:

\begin{verbatim}
Below are examples of correctly producing samples from the target
distribution. Each subsequent response should be drawn from the same
distribution.
\end{verbatim}

and the V6 system prompt is
\texttt{"When you respond to the next prompt, follow these steps:"}
followed by the task's reasoning lines (verbatim):

\begin{verbatim}
skewed_binary:
Step 1: Acknowledge that the target distribution is alpha with
probability 0.70 and beta with probability 0.30.
Step 2: Identify the support (alpha, beta).
Step 3: Mentally simulate a process that produces alpha with
probability 0.70 and beta with probability 0.30.
Step 4: Report the outcome of that simulation as your response.

bimodal_5way:
Step 1: Acknowledge that the target distribution is P(1)=0.40,
P(2)=0.10, P(3)=0.00, P(4)=0.10, P(5)=0.40.
Step 2: Identify the support (1, 2, 4, 5) - digit 3 has zero mass and
must never be produced.
Step 3: Mentally simulate a process that produces those frequencies.
Step 4: Report the outcome of that simulation as your response.

skewed_5way:
Step 1: Acknowledge that the target distribution is P(1)=0.50,
P(2)=0.20, P(3)=0.10, P(4)=0.10, P(5)=0.10.
Step 2: Identify the support (1, 2, 3, 4, 5).
Step 3: Mentally simulate a process that produces those frequencies.
Step 4: Report the outcome of that simulation as your response.
\end{verbatim}

(each \texttt{Step} is a single line in the code; the dash in the
bimodal\_5way Step 2 is an em-dash in the source).

% ------------------------------------------------------------
\subsection*{A.14\quad Stronger correction strategies V7--V9 (the stronger-corrections follow-up)}

Stronger correction variants (follow-up) on \texttt{uniform\_digit} and
\texttt{coin\_flip} (note this follow-up's coin task uses the user prompt
``Flip a fair coin. Report heads or tails. Reply with just the
word.'').

\textbf{V7 chain-of-thought with explicit probabilities} (system
prompt; user message is the task's base user prompt, 200 calls at
$T{=}1.0$):

\begin{verbatim}
uniform_digit:
For each request to pick a random digit from 1 to 5, think through the
following steps silently, then output only the final digit:
  Step 1: Acknowledge the target distribution. Uniform on {1,2,3,4,5}
means P(1)=0.2, P(2)=0.2, P(3)=0.2, P(4)=0.2, P(5)=0.2.
  Step 2: Compute the cumulative distribution: F(1)=0.2, F(2)=0.4,
F(3)=0.6, F(4)=0.8, F(5)=1.0.
  Step 3: Mentally draw a random number u in [0, 1].
  Step 4: Find the smallest digit d with F(d) >= u and output d.
Output only the final digit, nothing else.

coin_flip:
For each request to flip a fair coin, think through the following
steps silently, then output only the final result:
  Step 1: Acknowledge the target distribution. Fair coin:
P(heads)=0.5, P(tails)=0.5.
  Step 2: Mentally draw u in [0, 1].
  Step 3: If u < 0.5 output heads; otherwise output tails.
Output only the final result (heads or tails), nothing else.
\end{verbatim}

(each \texttt{Step} is a single line in the code).

\textbf{V8 list-of-N} (user message only, JSON response format, 10
calls of 20 samples each; \texttt{list\_size}=20):

\begin{verbatim}
uniform_digit:
Produce a list of 20 digits where the digits 1-5 each appear about 4
times (so the distribution across the 20 outputs is approximately
uniform on 1-5). Output the 20 digits as a JSON array under the key
'samples'.

coin_flip:
Produce a list of 20 coin-flip outcomes (each is either 'heads' or
'tails') where heads appears about 10 times and tails appears about 10
times (so the distribution is approximately fair). Output the 20
outcomes as a JSON array under the key 'samples'.
\end{verbatim}

\textbf{V9 sampler specification} (user message only, one call at
$T{=}0$, then 200 external NumPy draws, seed 42):

\begin{verbatim}
uniform_digit:
Output the probability distribution over digits 1 to 5 that represents
a uniform distribution. Output as JSON with key 'distribution' whose
value is a list of 5 probabilities [p1, p2, p3, p4, p5] summing to
1.0.

coin_flip:
Output the probability distribution over {heads, tails} that
represents a fair coin. Output as JSON with key 'distribution' whose
value is a list of 2 probabilities [p_heads, p_tails] summing to 1.0.
\end{verbatim}

% ------------------------------------------------------------
\subsection*{A.15\quad Alignment-ladder V0 tasks and base-model few-shot format}

Used by the post-training alignment-ladder experiment
(\texttt{silicon/ladder.py}: OLMo-2 base/SFT/DPO/Instruct,
Llama-3.1-8B base/instruct, Qwen2.5 base, Qwen2.5-Coder base,
Qwen2.5 instruct, and Qwen3-8B, run locally; V0 is 200 calls per task at $T{=}1.0$,
top-$p$ 1.0). The six tasks are the \texttt{point\_5} control plus the
five KNOWS/DOES tasks; the \texttt{bimodal\_5way} and
\texttt{skewed\_5way} targets follow the KNOWS/DOES panel (A.9) and
therefore differ from the seven-target panel's \texttt{bimodal\_15}
and \texttt{skewed\_5way} targets (A.8). The V0 prompts are:

\begin{verbatim}
point_5:       Always output the digit 5. Reply with just the digit.
uniform_digit: Pick a random digit from 1 to 5 uniformly at random.
               Reply with just the digit.
fair_coin:     Flip a fair coin. Output H or T.
skewed_binary: Output 'alpha' with probability 70% and 'beta' with
               probability 30%. Reply with just the word.
bimodal_5way:  Output a digit from 1 to 5 where 1 appears 40%, 2
               appears 10%, 3 appears 0%, 4 appears 10%, 5 appears
               40%. Reply with just the digit.
skewed_5way:   Output a digit from 1 to 5 where 1 appears 50%, 2
               appears 20%, 3 appears 10%, 4 appears 10%, 5 appears
               10%. Reply with just the digit.
\end{verbatim}

(each prompt is a single line in the code). Instruct-kind models
receive the V0 prompt zero-shot through their own chat template.
Base-kind models instead receive the fixed cross-domain 3-shot prefix
\texttt{BASE\_SHOTS} followed by \verb|Q: {v0_prompt}| and a newline
and \verb|A:|; the exemplars teach the bare-answer format without
demonstrating any answer from the test task's support:

\begin{verbatim}
Q: Pick a random letter from the English alphabet uniformly at random.
Output only the letter.
A: Q

Q: Pick a random integer uniformly between 1 and 100. Output only the
integer.
A: 37

Q: Pick a random color from red, green, blue. Output only the color.
A: green
\end{verbatim}

(each \texttt{Q:} line is a single line in the code). For the V8
list-of-20 condition (10 calls per task) the user message is built by
\texttt{v8\_user\_prompt}, the same construction as A.11 with
\texttt{\{freq\_str\}} enumerating expected counts
\texttt{round(p\_i * 20)} per support element:

\begin{verbatim}
Generate a list of 20 answers where {freq_str}. Output as JSON with
key 'samples' whose value is a list of 20 strings.
\end{verbatim}

Base-kind models receive the single canonical V8 exemplar
\texttt{BASE\_V8\_SHOT} followed by \verb|Q: {user}| and a newline and
\verb|A:|:

\begin{verbatim}
Q: Generate a list of 4 answers where 'X' about 2 times, 'Y' about 2
times. Output as JSON with key 'samples' whose value is a list of 4
strings.
A: {"samples": ["X", "Y", "Y", "X"]}
\end{verbatim}

(the \texttt{Q:} line is a single line in the code).

% ------------------------------------------------------------
\subsection*{A.16\quad Open-ended survey-domain tasks (four-model grid)}

Used by the open-ended perturbation experiment (\texttt{silicon/openended.py};
run with \texttt{python run.py openended --model <name>}). Each task has one
fixed prompt (below), six non-semantic whitespace/punctuation transforms of
that prompt applied programmatically, and ten semantic paraphrases drawn
uniformly per call; all paraphrases are reproduced verbatim here. The first entry of each list is the canonical prompt itself, so both perturbed conditions include the fixed prompt in their support.

\paragraph{important problem}
Fixed prompt:
\begin{verbatim}
In your view, what is the most important problem facing the United States today? Answer with a short phrase only.
\end{verbatim}
Paraphrases:
\begin{verbatim}
In your view, what is the most important problem facing the United States today? Answer with a short phrase only.
What do you consider the single biggest problem the country faces right now? Short phrase only.
Name the most pressing issue facing Americans today. Answer in a few words.
If you had to pick one, what is the top problem confronting the United States? Brief phrase only.
What national problem concerns you most at present? Answer with a short phrase.
Which issue facing the country today matters most? A few words only.
What is the most serious challenge the United States faces today? Short answer only.
Identify the most important issue for the nation right now. Phrase only.
In a national survey, you are asked for the country's most important problem. Your short answer:
What single problem should Americans be most worried about today? Brief phrase.
\end{verbatim}
\paragraph{news source}
Fixed prompt:
\begin{verbatim}
Which single news outlet do you rely on most for news? Output only the outlet's name.
\end{verbatim}
Paraphrases:
\begin{verbatim}
Which single news outlet do you rely on most for news? Output only the outlet's name.
What is your primary source of news? Name the outlet only.
Name the one news organization you turn to most often. Outlet name only.
Where do you get most of your news? Answer with a single outlet name.
If you could keep only one news source, which would it be? Name only.
Which news outlet do you check first each day? Output the name only.
What news source do you trust most? Outlet name only.
Name your main news outlet. Nothing but the name.
In a media survey, you are asked for your primary news source. Your answer (name only):
Which outlet provides most of the news you consume? Name only.
\end{verbatim}
\paragraph{policy priority}
Fixed prompt:
\begin{verbatim}
Name one policy area Congress should prioritize this year. Answer with a short phrase only.
\end{verbatim}
Paraphrases:
\begin{verbatim}
Name one policy area Congress should prioritize this year. Answer with a short phrase only.
Which policy issue should lawmakers focus on first this year? Short phrase only.
If Congress could advance one policy area this session, which should it be? Brief phrase.
What legislative priority matters most right now? Answer in a few words.
Name the policy domain most deserving of congressional attention. Phrase only.
Which area of policy should the federal government prioritize? Short answer only.
Pick one policy area for Congress to act on this year. A few words only.
What should top the legislative agenda this year? Brief phrase only.
In a civic survey, you are asked which policy area Congress should prioritize. Your short answer:
Which single policy issue should receive priority in Washington this year? Phrase only.
\end{verbatim}
\paragraph{respondent occupation}
Fixed prompt:
\begin{verbatim}
Generate a plausible occupation for a synthetic survey respondent. Output only the occupation.
\end{verbatim}
Paraphrases:
\begin{verbatim}
Generate a plausible occupation for a synthetic survey respondent. Output only the occupation.
Assign an occupation to a simulated survey participant. Occupation only.
A synthetic respondent needs an occupation field. Provide one. Output the occupation only.
For a simulated panelist profile, generate the occupation entry. Occupation only.
Produce one occupation for an artificial survey respondent. Output only the job title.
Fill in the occupation attribute for a generated respondent. Job title only.
What occupation should a synthetic panel member have? Output the occupation only.
Generate the employment field for a simulated respondent record. Occupation only.
Provide an occupation for a randomly generated survey persona. Occupation only.
Create one plausible occupation entry for a synthetic respondent. Output only the occupation.
\end{verbatim}
\paragraph{respondent city}
Fixed prompt:
\begin{verbatim}
Generate a plausible U.S. city of residence for a synthetic survey respondent. Output only the city name.
\end{verbatim}
Paraphrases:
\begin{verbatim}
Generate a plausible U.S. city of residence for a synthetic survey respondent. Output only the city name.
Assign a U.S. city of residence to a simulated survey participant. City name only.
A synthetic respondent needs a city field. Provide a U.S. city. Name only.
For a simulated panelist profile, generate the city of residence. U.S. city name only.
Produce one U.S. city for an artificial survey respondent's location. City only.
Fill in the residence city for a generated U.S. respondent. City name only.
In which U.S. city should a synthetic panel member live? Output the city only.
Generate the location field for a simulated respondent record. U.S. city only.
Provide a U.S. city of residence for a randomly generated survey persona. City only.
Create one plausible U.S. city entry for a synthetic respondent. Output only the city.
\end{verbatim}
